\documentclass[../../../include/open-logic-section]{subfiles}

\begin{document}

\olfileid[pt]{sth}{cardinals}{zfc}
\olsection{$\ZFC$: um marco}

Com a adição do axioma do bom ordenamento, atingimos o último
marco teórico. Dispomos agora de todos os axiomas necessários para~$\ZFC$.
Em detalhe:

\begin{defn}
A teoria $\ZFC$ tem estes axiomas: Extensionalidade, União, Pares,
Conjuntos potência, Infinitude, Fundação, Bom ordenamento e todas as instâncias dos esquemas de Separação e Substituição. Dito de outro modo, $\ZFC$ acrescenta
o Bom ordenamento a~$\ZF$.
\end{defn}

$\ZFC$ significa teoria de conjuntos de \emph{Zermelo--Fraenkel} com
\emph{Escolha}. Isto pode parecer ligeiramente estranho, pois o axioma
que adicionamos se chama ``Bom ordenamento'', não ``Escolha''. Mas, quando mais tarde
formularmos a \emph{Escolha}, verificar-se-á que o Bom ordenamento lhe é equivalente
(módulo~$\ZF$) (ver \olref[choice][woproblem]{thmwochoice}).
Logo, qual tomar como axioma ``básico'' é indiferente.
E o nome ``$\ZFC$'' é inteiramente standard na literatura.

\end{document}
